\documentclass{report}
\usepackage{amsmath,amssymb}
\setlength{\parindent}{0mm}
\setlength{\parskip}{1em}
\begin{document}
\begin{center}
\rule{6in}{1pt} \
{\large Edoardo Angelo Di Napoli \\
{\bf Block Iterative Solvers and Sequences of Eigenproblems arising in Electronic Structure Calculations}}

J\"ulich Supercomputing Centre \\ Institute for Advanced Simulation \\ Forschungszentrum J\"ulich GmbH \\ Wilhelm-Johnen Strasse \\ 52425 J\"ulich
\\
{\tt e.di.napoli@fz-juelich.de}\end{center}

Research in several branches of chemistry and materials science
relies on large ab initio numerical simulations. The majority of these
simulations are based on computational methods developed within the
framework of Density Functional Theory (DFT). DFT provides the means to
solve a high-dimensional quantum mechanical problem by transforming it
into a large set of coupled one-dimensional equations, which is
ultimately represented as a non-linear generalized eigenvalue problem.
The latter is solved self-consistently through a series of successive
iteration cycles: the solution computed at the end of one cycle is used
to generate the input in the next until the distance between two
successive solutions is negligible.

Typically a simulations requires tens of cycles before reaching
convergence. After the discretization -- intended in the general sense of
reducing a continuous problem to one with a finite number of unknowns --
each cycle comprises dozens of large generalized eigenproblems
$P^{(i)}_{\bf k}: A^{(i)}_{\bf k} x - \lambda B^{(i)}_{\bf k} x$ where
$A$ is hermitian and $B$ hermitian positive definite. Within every cycle
the eigenproblems are parametrized by the reciprocal lattice vector ${\bf
k}$ while the iteration index $i$ denotes the iteration cycle. The size
of each problem ranges from 10 to 40 thousand and the interest lays in
the eigenpairs corresponding to the lower 10-20\% part of the spectrum.
Due to the dense nature of the eigenproblems and the large portion of the
spectrum requested, iterative solvers are generally not competitive; as a
consequence, current simulation codes uniquely use direct methods.

Generally, the eigenproblem $P^{(i+1)}_{\bf k}$ is solved in complete
isolation from all $P^{(i)}$s. This is surprising in light of the fact
that each $P^{(i+1)}_{\bf k}$ is constructed by manipulating the
solutions of all the problems at iteration $i$. In a recent study, the
author harnessed this last observation by considering every simulation as
made of dozens of sequences $\{P_{\bf k}\}$, where each sequence groups
together eigenproblems with equal {\bf k}-vector and increasing iteration
index $i$. It was then demonstrated that successive eigenproblems in a
sequence are strongly correlated to one another. In particular, by
tracking the evolution over iterations of the angle between eigenvectors
$x^{(i)}$ and $x^{(i+1)}$, it was shown the angles decrease noticeably
after the first few iterations. Even though the overall simulation has
not yet converged, the eigenvectors of $P^{(i)}_{\bf k}$ and
$P^{(i+1)}_{\bf k}$ are close to collinear.

This result suggests we could use the eigenvectors of $P^{(i)}$ as an
educated guess for the eigenvectors of the successive problem
$P^{(i+1)}$. The key element is to exploit the collinearity between
vectors of adjacent problems in order to improve the performance of the
solver. While implementing this strategy naturally leads to the use of
iterative solvers, not all the methods are well suited to the task at
hand.

In light of these considerations, exploiting the evolution of
eigenvectors depends on two distinct key properties of the iterative
method of choice: 1) the ability to receive as input a sizable set of
approximate eigenvectors and exploit them to efficiently compute the
required eigenpairs, and 2) the capacity to solve simultaneously for a
substantial portion of eigenpairs. The first property implies the solver
achieves a moderate-to-substantial reduction in the number of
matrix-vector operations as well as an improved convergence. The second
characteristic results in a solver capable of filtering away as
efficiently as possible the unwanted components of the approximate
eigenvectors.

In accord to these requirements, the class of block iterative methods
constitutes the natural choice in order to satisfy property 1); by
accepting a variable set of multiple starting vectors, these methods have
a faster convergence rate and avoid stalling when facing small clusters
of eigenvalues. When block methods are augmented with polynomial
accelerators they also satisfy property 2) and their performance is
further improved.

In this study we present preliminary results that would eventually open
the way to a widespread use of iterative solvers in ab initio electronic
structure codes. We provide numerical examples where opportunely selected
iterative solvers benefit from the re-use of eigenvectors when applied to
sequences of eigenproblems extracted from simulations of real-world
materials. At first we experimented with a block method (block
Krylov-Schur) satisfying only property 1). At a later stage we selected a
solver satisfying both properties (block Chebyshev-Davidson) and
performed similar experiments.

In our investigation we vary several parameters in order to address how
the solvers behave under different conditions. We selected different size
for the sequence of eigenproblems as well as an increasing number of
eigenpairs to be computed. We also vary the block size in relation with
the fraction of eigenpairs sought after and analyze its influence on the
performance of the solver. In most cases our study shows that, when the
solvers are fed approximated eigenvectors as opposed to random vectors,
they obtain a substantial speed-up. The increase in performance changes
with the variation of the parameters but strongly indicates that
iterative solvers could become a valid alternative to direct methods.

The pivotal element allowing the achievement of such a result resides in
the transposition of the ``matrix'' of eigenproblems emerging from
electronic structure calculations; while the computational physicist
solves many rows of {\bf k}-eigenproblems (one for each cycle) we look at
the entire computation as made of {\bf k}-sequences. This shift in focus
enables one to harness the inherent essence of the iterative cycles and
translate it to an efficient use of opportunely tailored iterative
solvers. In the long term such a result will positively impact a large
community of computational scientists working on DFT-based methods.



\end{document}
